\scanpage{12}
\[
\sigma_\infty(g^{-1})=l_0^\beta g l_1^\beta.
\]
Which determines $\beta$ uniquely as a function of $g$.

\deleted{\unsure{Applying|Let us apply} the fundamental relation}
\[
\deleted{\displaystyle u\unsure{g|\rho}=\operatorname{int}(g)\rho u}
\]
\deleted{to $L_0$, we find, since $u(L_0)=L_0$, $\rho(L_0)=L_1$}
\[
\deleted{\displaystyle u(L_1)=\operatorname{int}(g)(L_{\unsure{1|0}})}
\]
\deleted{and this already determines \deleted{\unsure{$\operatorname{int}(g)$}} \unsure{$g$|\deleted{$g$}} modulo right multiplication by an element of $L_1$, \unsure{by|near|then} the datum of $u$}

We have prettier generators of the \inserted{sub-}extension of $\unsure{\mathfrak S_3|\mathcal S_3}$ by $\unsure{\pi|\Pi}$ \inserted{in $\operatorname{Aut}_{\mathrm{lac}}(\unsure{\pi|\Pi})$,} than the six generators $l_0,l_1,l_\infty,\sigma_0,\sigma_1,\sigma_\infty$ (giving $\rho$ by $\sigma_0\sigma_1=l_1^{-1}\rho$ \deleted{\unsure{for example}} i.e. $\rho=\sigma_0\sigma_1l_1$ for example), by introducing the ramification groups at the points $0,1,\infty$, which are canonicall\textsuperscript{y} isomorphic \unsure{to|:} $\widehat{\mathbb Z}$,
We shall therefore define $\varepsilon_0,\varepsilon_1,\varepsilon_\infty$ by the relations
\[
\varepsilon_0^2=l_0,\qquad\varepsilon_1^2=l_1,\qquad\varepsilon_\infty^2=l_\infty
\]
and more precisely, $\varepsilon_i$ will be the unique \deleted{\unsure{elemen\textsuperscript{t}}} \inserted{automorphism} in the class modulo $\unsure{\pi|\Pi}$ of the $\sigma_i$, \deleted{and} which fixes $l_i$ and is \unsure{\deleted{\unsure{of|$\sigma$}}|$\sigma$} \inserted{\unsure{of square $l_i$}}
\mspage{11}
